\begin{pastexam}[label={px:gw3-112-ds-04}]{關務三等資訊處理 112 年 第四題}
（申論；25 分）

給予如下程式片段，假設 x[] = [25, 57, 48, 37, 12, 92, 86, 33]，請只用下述 C 語言宣告的變數及兩個 for 迴圈，完成下面的
選擇排序（Selection Sort），假設有 n 個資料要由小排到大，每一外迴圈將最大值放在第 n-1 個位置，然後第二大的資料放在第
n-2 個位置，依此類推，將資料放到適當的位置，執行後陣列 x[]內容由小排至大。（25 分）
\begin{minted}[fontsize=\small]{c}
Selectsort(x, n)
int x[], n;
{
    int i, index, j, large;

    for (i = n-1; i > 0; i--){

        for (j = 1; j <= i; j++){

        }
    }
}
\end{minted}
\end{pastexam}

\begin{pastsoln}
答案：外迴圈開始時令 \texttt{large = x[0]; index = 0;}，內迴圈遇到 \texttt{x[j] > large} 就更新 \texttt{large}、\texttt{index}，
內迴圈結束後 \texttt{x[index] = x[i]; x[i] = large;}。 排序結果 $12, 25, 33, 37, 48, 57, 86, 92$。

\textbf{想法}（資結書第~14.2~節的選擇排序，這裡是「每次選最大的放到後面」的版本）：第 $i$ 輪（$i = n-1, n-2, \ldots, 1$）
在 $x[0..i]$ 中找最大值，記下它的值 \texttt{large} 和位置 \texttt{index}，再和 $x[i]$ 交換，$x[i..n-1]$ 就排好了。 內迴圈從
$j = 1$ 開始，所以要先拿 $x[0]$ 當目前的最大值。 交換時 \texttt{large} 本身就存著最大值，不需要另外的暫存變數，正好只用
題目宣告的四個變數。
\begin{minted}[fontsize=\small]{c}
Selectsort(x, n)
int x[], n;
{
    int i, index, j, large;

    for (i = n-1; i > 0; i--){
        large = x[0];              /* 先假設 x[0] 是 x[0..i] 的最大值 */
        index = 0;
        for (j = 1; j <= i; j++){
            if (x[j] > large) {    /* 找到更大的，記下值和位置 */
                large = x[j];
                index = j;
            }
        }
        x[index] = x[i];           /* 最大值和 x[i] 交換：x[i] 搬到最大值原來的位置 */
        x[i] = large;              /* 最大值放到 x[i] */
    }
}
\end{minted}
（\texttt{Selectsort(x, n)} 加上 \texttt{int x[], n;} 是舊式的 K\&R 函式定義，沒寫回傳型別時預設為 \texttt{int}；現代 C 寫成
\texttt{void Selectsort(int x[], int n)}。）

\textbf{追蹤} $x = [25, 57, 48, 37, 12, 92, 86, 33]$，$n = 8$：
\begin{center}
\small
\begin{tabular}{cccl}
\toprule
$i$ & 最大值 large & index & 交換後的 $x[0..7]$ \\
\midrule
7 & 92 & 5 & 25, 57, 48, 37, 12, 33, 86, \textbf{92} \\
6 & 86 & 6 & 25, 57, 48, 37, 12, 33, \textbf{86}, \textbf{92} \\
5 & 57 & 1 & 25, 33, 48, 37, 12, \textbf{57}, \textbf{86}, \textbf{92} \\
4 & 48 & 2 & 25, 33, 12, 37, \textbf{48}, \textbf{57}, \textbf{86}, \textbf{92} \\
3 & 37 & 3 & 25, 33, 12, \textbf{37}, \textbf{48}, \textbf{57}, \textbf{86}, \textbf{92} \\
2 & 33 & 1 & 25, 12, \textbf{33}, \textbf{37}, \textbf{48}, \textbf{57}, \textbf{86}, \textbf{92} \\
1 & 25 & 0 & \textbf{12}, \textbf{25}, \textbf{33}, \textbf{37}, \textbf{48}, \textbf{57}, \textbf{86}, \textbf{92} \\
\bottomrule
\end{tabular}
\end{center}
（粗體是已經排好的部分。 $i = 6$、$i = 3$ 時最大值本來就在 $x[i]$，「交換」只是寫回原位，不影響結果。）最後
$x = [12, 25, 33, 37, 48, 57, 86, 92]$，由小到大。

\textbf{正確性與複雜度}：每輪結束後 $x[i]$ 是 $x[0..i]$ 的最大值，而 $x[i+1..n-1]$ 是更早幾輪放好的、比它們都大的值，
所以 $x[i..n-1]$ 已排好；$i = 1$ 那輪結束後 $x[0] \le x[1]$，全部排好。 內迴圈比較 $i$ 次，總比較次數
$\sum_{i=1}^{n-1} i = n(n-1)/2$，本例 $28$ 次；時間 $\Theta(n^2)$，交換最多 $n - 1$ 次，額外空間 $O(1)$。
\end{pastsoln}
